\begin{table}[!htbp]
\centering
\small
\caption{School-type-specific total VA slopes in submitted rankings}
\label{tab:school-type-total-va-slopes}
\begin{tabular}{llcc}
\toprule
School type & VA dimension & Published coordinates & Comuna centers \\
\midrule
Non-TP & Math VA & 0.0510 (0.0021) & 0.0502 (0.0021) \\
Non-TP & Exam-taking VA & 0.1535 (0.0043) & 0.1461 (0.0043) \\
Non-TP & High-premium-field VA & 0.0114 (0.0031) & 0.0091 (0.0031) \\
Non-TP & High-premium-institution VA & -0.0083 (0.0031) & -0.0136 (0.0031) \\
\midrule
Mixed HC/TP & Math VA & 0.0576 (0.0050) & 0.0546 (0.0049) \\
Mixed HC/TP & Exam-taking VA & 0.0729 (0.0032) & 0.0729 (0.0032) \\
Mixed HC/TP & High-premium-field VA & 0.0626 (0.0021) & 0.0598 (0.0021) \\
Mixed HC/TP & High-premium-institution VA & -0.0573 (0.0041) & -0.0613 (0.0041) \\
\midrule
TP only & Math VA & -0.0063 (0.0087) & -0.0161 (0.0086) \\
TP only & Exam-taking VA & 0.0256 (0.0047) & 0.0348 (0.0047) \\
TP only & High-premium-field VA & 0.0578 (0.0024) & 0.0578 (0.0023) \\
TP only & High-premium-institution VA & -0.0229 (0.0044) & -0.0237 (0.0044) \\
\bottomrule
\end{tabular}
\par\smallskip
\begin{minipage}{0.98\linewidth}\footnotesize
Each cell reports a total VA slope with its student-clustered standard error. Non-TP slopes equal the base VA coefficient. Mixed HC/TP and TP-only slopes equal the base plus the corresponding interaction. SEs use the exact covariance of the coefficient contrast. All VA indices use the same pooled school normalization.
\end{minipage}
\end{table}
